\documentclass[12pt,a4paper]{article}

% ------------------------------------------------------------
% PACKAGES
% ------------------------------------------------------------
\usepackage[margin=1in]{geometry}
\usepackage[T1]{fontenc}
\usepackage[utf8]{inputenc}
\usepackage{lmodern}
\usepackage{enumitem}
\usepackage{xcolor}
\usepackage{hyperref}
\usepackage{array}
\usepackage{tabularx}
\usepackage{booktabs}
\usepackage{longtable}

\hypersetup{
	colorlinks=true,
	linkcolor=black,
	urlcolor=blue
}

\setcounter{secnumdepth}{0}

% ------------------------------------------------------------
% TITLE
% ------------------------------------------------------------
\title{Tableau Visualization Strategy\\
	for the Formula 1 Data Warehouse Project}
\author{Giuseppe Rudi}
\date{\today}

\begin{document}
	
	\maketitle
	
	\tableofcontents
	\newpage
	
	% ------------------------------------------------------------
	\section{Purpose of this Document}
	
	This document defines the planned structure of the visualization phase of the Formula 1 Data Warehouse project.
	
	The goal is not to describe SQL queries or final technical implementation details.
	At this stage, the objective is to define the analytical storytelling, the number of dashboards, the purpose of each dashboard, the types of charts to build in Tableau, and the logical connection between them.
	
	The visualization must support a clear five-minute presentation.
	For this reason, the dashboard system must be compact, visually coherent, and strongly connected to the main Formula 1 story.
	
	The final visualization must not be a collection of independent charts.
	It must work as a guided analytical story where each dashboard answers a specific question, and all the answers together support the final interpretation.
	
	% ------------------------------------------------------------
	\section{Main Storytelling}
	
	The central question of the visualization is:
	
	\begin{quote}
		\textit{How did the competitive balance among Ferrari, Red Bull, and Mercedes change from the 2021 season to the 2022 season after the regulatory change?}
	\end{quote}
	
	The story does not simply ask which team was faster.
	It tries to understand:
	
	\begin{itemize}[leftmargin=1.2cm]
		\item whether the competitive balance changed between 2021 and 2022;
		\item whether the change was the same in Qualifying and Race sessions;
		\item whether the change depended on circuit characteristics;
		\item where inside the lap the performance differences appeared;
		\item whether tyre compounds and straight-line speed indicators help explain the observed differences.
	\end{itemize}
	
	The final message of the visualization can be summarized as follows:
	
	\begin{quote}
		\textit{The 2022 regulation change did not only change which team was faster. It changed how, where, and under which technical conditions Ferrari, Red Bull, and Mercedes were competitive.}
	\end{quote}
	
	% ------------------------------------------------------------
	\section{Final Dashboard Structure}
	
	For the final presentation, the selected structure contains four dashboards.
	
	\begin{enumerate}[leftmargin=1.2cm]
		\item \textbf{Competitive Balance Overview}
		\item \textbf{Qualifying vs Race Performance}
		\item \textbf{Circuit and Sector DNA}
		\item \textbf{Tyre and Straight-Line Performance}
	\end{enumerate}
	
	This structure is preferred because it is detailed enough to tell a complete story, but compact enough to be presented in five minutes.
	
	A larger number of dashboards would create the risk of moving too quickly and not explaining the analytical meaning of each chart.
	A smaller number of dashboards would force too many different ideas into the same page, making the final dashboard less readable.
	
	% ------------------------------------------------------------
	\section{Why Four Dashboards Are the Best Choice}
	
	The project has a strong Formula 1 story, but the presentation time is limited.
	For this reason, the dashboard sequence must follow a clear progression:
	
	\begin{center}
		\begin{tabularx}{\textwidth}{p{4cm} X}
			\toprule
			\textbf{Dashboard} & \textbf{Narrative Role} \\
			\midrule
			Competitive Balance Overview & Shows what changed between 2021 and 2022. \\
			Qualifying vs Race Performance & Explains whether the change depends on the session context. \\
			Circuit and Sector DNA & Explains where the advantage appears, using circuit and sector characteristics. \\
			Tyre and Straight-Line Performance & Adds a technical interpretation based on tyre compounds and speed indicators. \\
			\bottomrule
		\end{tabularx}
	\end{center}
	
	The logic is:
	
	\begin{quote}
		First we show the global change, then we separate the performance context, then we explain the circuit and sector behaviour, and finally we add tyres and straight-line speed as technical interpretation factors.
	\end{quote}
	
	% ------------------------------------------------------------
	\section{Why Wet Performance Is Not a Main Dashboard}
	
	Wet performance is not selected as one of the four main dashboards.
	
	This does not mean that wet sessions are useless or invalid.
	A wet lap is a valid Formula 1 observation.
	However, wet conditions introduce many additional variables, such as:
	
	\begin{itemize}[leftmargin=1.2cm]
		\item changing track grip;
		\item different tyre compounds;
		\item different driver confidence;
		\item changing weather intensity;
		\item different setup compromises;
		\item Safety Car or red flag probability;
		\item traffic and visibility effects.
	\end{itemize}
	
	For this reason, wet performance is less suitable for the main comparison between Ferrari, Red Bull, and Mercedes.
	
	The main storytelling focuses on clean and dry conditions because this produces a more comparable analytical subset.
	Wet performance can be mentioned as a possible future extension or as an optional filter, but it should not become one of the main dashboards.
	
	Suggested sentence for the report:
	
	\begin{quote}
		Wet sessions are excluded from the main storytelling because they introduce too many additional factors and would make the comparison less controlled. They are valid data, but they are outside the analytical focus of this visualization.
	\end{quote}
	
	% ------------------------------------------------------------
	\section{Why We Do Not Call the Fourth Dashboard ``Engine Performance''}
	
	The fourth dashboard includes speed-related metrics.
	However, it should not be called \textit{Engine Performance}.
	
	The reason is that the Data Warehouse does not contain direct information about the internal power unit, engine modes, deployment strategy, aerodynamic drag, wing configuration, or team setup decisions.
	
	Speed values observed on a lap are influenced by many factors:
	
	\begin{itemize}[leftmargin=1.2cm]
		\item power unit performance;
		\item aerodynamic drag;
		\item rear and front wing setup;
		\item DRS availability;
		\item slipstream;
		\item tyre compound and tyre condition;
		\item fuel load;
		\item wind direction and wind speed;
		\item circuit layout;
		\item traffic and race context.
	\end{itemize}
	
	Therefore, the correct interpretation is not:
	
	\begin{quote}
		\textit{This team has the best engine.}
	\end{quote}
	
	The correct interpretation is:
	
	\begin{quote}
		\textit{This team shows stronger straight-line speed indicators in the selected context.}
	\end{quote}
	
	For this reason, the fourth dashboard should be called:
	
	\begin{quote}
		\textbf{Tyre and Straight-Line Performance}
	\end{quote}
	
	This title is more precise and more defensible during the exam.
	
	% ------------------------------------------------------------
	\section{General Visual Style}
	
	The dashboard should visually recall the Formula 1 world.
	
	The objective is to create a dashboard that looks modern, technical, and close to a race-control or F1 broadcast style.
	
	\subsection{Recommended Visual Identity}
	
	\begin{itemize}[leftmargin=1.2cm]
		\item dark background;
		\item white or light grey text;
		\item red accent color;
		\item team colors for Ferrari, Red Bull, and Mercedes;
		\item compact KPI cards;
		\item minimal grid lines;
		\item strong contrast;
		\item clean spacing;
		\item short explanatory text;
		\item clear legends.
	\end{itemize}
	
	\subsection{Recommended Font Style}
	
	The dashboard should use a technical and modern font.
	
	If the official Formula 1 font is not legally available, it is better to use a similar but standard font.
	
	Possible choices are:
	
	\begin{itemize}[leftmargin=1.2cm]
		\item Bahnschrift;
		\item Rajdhani;
		\item Orbitron;
		\item DIN-style fonts;
		\item Eurostile-like fonts.
	\end{itemize}
	
	The objective is not to copy the official Formula 1 identity illegally, but to create a coherent visual style inspired by Formula 1 graphics.
	
	% ------------------------------------------------------------
	\section{Common Dashboard Structure}
	
	All dashboards should follow the same structure.
	
	\begin{verbatim}
		------------------------------------------------------------
		| Dashboard Title / Main Question                          |
		| Short subtitle explaining the purpose                    |
		------------------------------------------------------------
		| Global filters                                           |
		------------------------------------------------------------
		| KPI cards or summary strip                               |
		------------------------------------------------------------
		| Main chart                                               |
		------------------------------------------------------------
		| Supporting chart 1             | Supporting chart 2       |
		------------------------------------------------------------
		| Short interpretation / key message                       |
		------------------------------------------------------------
	\end{verbatim}
	
	This repeated structure is important because it gives coherence to the entire visualization.
	The user should not feel that each dashboard belongs to a different project.
	
	% ------------------------------------------------------------
	\section{Common Filters}
	
	The dashboards should share a common set of filters whenever possible.
	
	\begin{longtable}{p{4cm} p{10cm}}
		\toprule
		\textbf{Filter} & \textbf{Purpose} \\
		\midrule
		Season & Allows comparison between 2021 and 2022. \\
		Team & Allows focus on Ferrari, Red Bull, Mercedes, or a subset of them. \\
		Session Type & Allows comparison between Qualifying and Race. \\
		Grand Prix & Allows selection of a specific race weekend. \\
		Circuit Category & Allows analysis by track profile. \\
		Tyre Compound & Allows focus on Soft, Medium, and Hard tyres. \\
		Clean Lap Flag & Allows the analysis to focus on comparable laps only. \\
		\bottomrule
	\end{longtable}
	
	The default configuration should show the comparison between 2021 and 2022 together.
	The user should not be forced to select 2021 first and then 2022 manually.
	Since the main question is comparative, the dashboard should already show the two seasons side by side by default.
	
	% ------------------------------------------------------------
	\section{Dashboard 1: Competitive Balance Overview}
	
	\subsection{Main Question}
	
	\begin{quote}
		\textit{How did the competitive balance among Ferrari, Red Bull, and Mercedes change from 2021 to 2022?}
	\end{quote}
	
	\subsection{Purpose}
	
	This is the opening dashboard.
	It introduces the entire story and answers the first question: \textit{what changed?}
	
	The user must immediately understand:
	
	\begin{itemize}[leftmargin=1.2cm]
		\item which teams are compared;
		\item which seasons are compared;
		\item whether the competitive gap changed;
		\item whether one team improved or declined relative to the others.
	\end{itemize}
	
	\subsection{Recommended Charts}
	
	\begin{longtable}{p{4cm} p{4cm} p{6cm}}
		\toprule
		\textbf{Chart} & \textbf{Type} & \textbf{Purpose} \\
		\midrule
		KPI Cards & Summary cards & Show selected laps, selected Grand Prix, best team, and average performance gap. \\
		Season Comparison & Slope chart or grouped bar chart & Compare team performance between 2021 and 2022. \\
		Gap Trend by Grand Prix & Line chart & Show how the gap evolves across the season. \\
		Team Ranking Matrix & Heatmap & Show relative strength by season and session context. \\
		\bottomrule
	\end{longtable}
	
	\subsection{Why These Charts Are Chosen}
	
	The KPI cards create an immediate overview.
	The slope chart is useful because the main comparison is between two seasons.
	The line chart shows whether the difference is constant or changes across Grand Prix.
	The heatmap gives a compact view of relative performance.
	
	\subsection{Dashboard Message}
	
	Suggested text to include inside Tableau:
	
	\begin{quote}
		This dashboard provides the global view of the competitive balance between Ferrari, Red Bull, and Mercedes across the 2021 and 2022 seasons. The comparison is based on clean and comparable laps.
	\end{quote}
	
	\subsection{Transition to the Next Dashboard}
	
	\begin{quote}
		After observing the global competitive change, the next step is to understand whether this change appears in the same way in Qualifying and Race sessions.
	\end{quote}
	
	% ------------------------------------------------------------
	\section{Dashboard 2: Qualifying vs Race Performance}
	
	\subsection{Main Question}
	
	\begin{quote}
		\textit{Was the performance change the same in Qualifying and Race sessions?}
	\end{quote}
	
	\subsection{Purpose}
	
	This dashboard separates two different performance contexts:
	
	\begin{itemize}[leftmargin=1.2cm]
		\item Qualifying, which approximates one-lap pace;
		\item Race, which represents a more complex performance context.
	\end{itemize}
	
	This distinction is important because a Formula 1 car can be strong in Qualifying but weaker in Race conditions, or vice versa.
	
	\subsection{Recommended Charts}
	
	\begin{longtable}{p{4cm} p{4cm} p{6cm}}
		\toprule
		\textbf{Chart} & \textbf{Type} & \textbf{Purpose} \\
		\midrule
		Qualifying Gap by Team & Bar chart & Compare one-lap pace between teams. \\
		Race Gap by Team & Bar chart & Compare race pace between teams. \\
		Qualifying vs Race Comparison & Scatterplot & Show whether a team is balanced or stronger in one context. \\
		Session Trend & Line chart & Show performance evolution by session type. \\
		\bottomrule
	\end{longtable}
	
	\subsection{Why These Charts Are Chosen}
	
	The two bar charts create a clear comparison between Qualifying and Race.
	The scatterplot helps identify whether a team is strong in both contexts or only in one.
	The line chart shows whether this difference changes across Grand Prix.
	
	\subsection{Dashboard Message}
	
	Suggested text to include inside Tableau:
	
	\begin{quote}
		This dashboard separates one-lap pace from race pace. The goal is to verify whether the competitive change between 2021 and 2022 is visible in both session types or only in a specific performance context.
	\end{quote}
	
	\subsection{Transition to the Next Dashboard}
	
	\begin{quote}
		After separating Qualifying and Race performance, the next step is to understand whether the differences depend on the technical characteristics of circuits and sectors.
	\end{quote}
	
	% ------------------------------------------------------------
	\section{Dashboard 3: Circuit and Sector DNA}
	
	\subsection{Main Question}
	
	\begin{quote}
		\textit{Where was the lap time advantage created?}
	\end{quote}
	
	\subsection{Purpose}
	
	This dashboard is the most technical and distinctive part of the storytelling.
	
	It connects team performance with:
	
	\begin{itemize}[leftmargin=1.2cm]
		\item global circuit category;
		\item sector category;
		\item sector-level lap time;
		\item circuit-specific behaviour.
	\end{itemize}
	
	The goal is to move from a simple comparison of lap times to a more Formula 1-oriented interpretation.
	
	\subsection{Interpretation Logic}
	
	Each circuit is classified according to its global technical profile.
	Each sector is also classified according to its dominant characteristic.
	
	For example:
	
	\begin{itemize}[leftmargin=1.2cm]
		\item a \texttt{Power} sector is dominated by acceleration zones and long straights;
		\item a \texttt{FastCorners} sector is useful to observe aerodynamic stability;
		\item a \texttt{SlowCorners} sector may highlight traction and mechanical balance;
		\item a \texttt{Technical} sector may highlight precision, direction changes, and car balance.
	\end{itemize}
	
	This does not allow us to measure aerodynamic load or engine performance directly.
	However, it allows us to observe performance patterns that are compatible with specific technical contexts.
	
	\subsection{Recommended Charts}
	
	\begin{longtable}{p{4cm} p{4cm} p{6cm}}
		\toprule
		\textbf{Chart} & \textbf{Type} & \textbf{Purpose} \\
		\midrule
		Team by Circuit Category & Heatmap & Show which team performs better on Street, PowerSensitive, AeroSensitive, and Mixed circuits. \\
		Performance by Circuit Category & Grouped bar chart & Compare average gap by circuit type. \\
		Sector Gap by Team & Bar chart & Show where teams gain or lose time in Sector 1, Sector 2, and Sector 3. \\
		Circuit Sector Profile & KPI/text panel & Show the selected circuit and its sector categories. \\
		\bottomrule
	\end{longtable}
	
	\subsection{Why These Charts Are Chosen}
	
	The heatmap is useful because it compares teams and circuit categories at the same time.
	The grouped bar chart provides a more readable comparison by category.
	The sector gap chart decomposes the lap time advantage.
	The circuit profile card gives context and makes the dashboard easier to interpret.
	
	\subsection{Dashboard Message}
	
	Suggested text to include inside Tableau:
	
	\begin{quote}
		This dashboard links lap performance with the technical identity of each circuit and sector. The objective is to understand where the advantage appears, without claiming to directly measure hidden technical parameters such as aerodynamic load or engine maps.
	\end{quote}
	
	\subsection{Important Analytical Limitation}
	
	This dashboard must be interpreted carefully.
	
	Formula 1 performance depends on many factors that are not fully visible in public data:
	
	\begin{itemize}[leftmargin=1.2cm]
		\item aerodynamic setup;
		\item wing configuration;
		\item power unit mode;
		\item tyre condition;
		\item fuel load;
		\item traffic;
		\item DRS;
		\item slipstream;
		\item weather and wind;
		\item team-specific setup choices.
	\end{itemize}
	
	Therefore, the dashboard does not prove absolute technical superiority.
	It shows performance evidence under selected assumptions.
	
	Suggested report sentence:
	
	\begin{quote}
		The dashboard provides evidence of performance patterns in different circuit and sector contexts, but it does not claim to isolate the exact technical cause of those patterns.
	\end{quote}
	
	\subsection{Transition to the Next Dashboard}
	
	\begin{quote}
		After identifying where the advantage appears on circuits and sectors, the final dashboard adds two important interpretation factors: tyre compound and straight-line speed indicators.
	\end{quote}
	
	% ------------------------------------------------------------
	\section{Dashboard 4: Tyre and Straight-Line Performance}
	
	\subsection{Main Question}
	
	\begin{quote}
		\textit{How do tyre compound and straight-line speed indicators influence the interpretation of team performance?}
	\end{quote}
	
	\subsection{Purpose}
	
	This dashboard completes the story by adding two technical layers:
	
	\begin{itemize}[leftmargin=1.2cm]
		\item tyre compound behaviour;
		\item speed indicators in power-sensitive contexts.
	\end{itemize}
	
	The objective is to understand whether some teams appear stronger with specific tyre compounds and whether speed metrics highlight differences in sectors or circuits where straight-line performance matters.
	
	\subsection{Why Tyres Are Included}
	
	Tyres are central in Formula 1 performance.
	A car can be strong on Soft tyres but less consistent on Medium or Hard tyres.
	Another car may be less explosive on Soft but more stable on harder compounds.
	
	For this reason, tyre compound analysis helps explain performance differences that may not be visible from average lap time alone.
	
	\subsection{Why Straight-Line Speed Is Included}
	
	The lap data include speed-related attributes.
	These attributes can be used to study straight-line speed indicators, especially in power-sensitive circuits or sectors.
	
	However, the dashboard must not claim to measure engine performance directly.
	Speed is influenced by many factors, including power unit, drag, setup, DRS, slipstream, tyre condition, fuel load, and weather.
	
	For this reason, the correct interpretation is:
	
	\begin{quote}
		\textit{The dashboard observes straight-line speed indicators, not pure engine performance.}
	\end{quote}
	
	\subsection{Recommended Charts}
	
	\begin{longtable}{p{4cm} p{4cm} p{6cm}}
		\toprule
		\textbf{Chart} & \textbf{Type} & \textbf{Purpose} \\
		\midrule
		Compound Performance Distribution & Box plot & Show lap time or gap distribution by tyre compound. \\
		Team by Compound & Heatmap & Identify which team performs better on Soft, Medium, and Hard. \\
		Speed Metrics by Team & Bar chart & Compare speed indicators such as speed trap or maximum speed values. \\
		Speed in Power Contexts & Scatterplot or bar chart & Compare speed metrics only on PowerSensitive circuits or Power sectors. \\
		\bottomrule
	\end{longtable}
	
	\subsection{Why These Charts Are Chosen}
	
	The box plot is useful because tyre performance is not only about the average.
	It also shows variability and consistency.
	The heatmap gives a compact view of team-compound behaviour.
	The speed chart adds a technical layer.
	The power-context chart avoids mixing circuits where straight-line speed is less meaningful.
	
	\subsection{Dashboard Message}
	
	Suggested text to include inside Tableau:
	
	\begin{quote}
		This dashboard analyzes tyre compound behaviour and straight-line speed indicators. Speed values are interpreted as public-data proxies and not as direct measurements of power unit performance.
	\end{quote}
	
	\subsection{Important Analytical Limitation}
	
	Suggested report sentence:
	
	\begin{quote}
		Straight-line speed is not interpreted as pure engine performance, because it is affected by aerodynamic drag, setup, DRS, slipstream, tyre condition, fuel load, wind, and race context. Therefore, the dashboard uses speed metrics only as indicators of straight-line behaviour in selected contexts.
	\end{quote}
	
	% ------------------------------------------------------------
	\section{Five-Minute Presentation Plan}
	
	The dashboard presentation must be designed for a five-minute pitch.
	
	A possible structure is:
	
	\begin{longtable}{p{3cm} p{5cm} p{6cm}}
		\toprule
		\textbf{Time} & \textbf{Dashboard} & \textbf{Message} \\
		\midrule
		0:00--0:30 & Introduction & Present the main question and the 2021--2022 regulatory context. \\
		0:30--1:40 & Competitive Balance Overview & Show the global change in team performance. \\
		1:40--2:40 & Qualifying vs Race & Explain whether the change is the same in one-lap pace and race pace. \\
		2:40--4:00 & Circuit and Sector DNA & Show where the advantage appears across circuit and sector types. \\
		4:00--4:40 & Tyre and Straight-Line Performance & Add tyre and speed interpretation, with clear limitations. \\
		4:40--5:00 & Conclusion & Summarize the final answer to the main question. \\
		\bottomrule
	\end{longtable}
	
	The presentation should not spend too much time on technical implementation.
	The focus must remain on the story and on the interpretation of the dashboards.
	
	% ------------------------------------------------------------
	\section{Suggested Final Oral Conclusion}
	
	A possible final explanation is:
	
	\begin{quote}
		The dashboard system shows that the 2022 regulation change should not be interpreted only as a change in ranking. The competitive balance between Ferrari, Red Bull, and Mercedes must be analyzed across different contexts: season, session type, circuit category, sector type, tyre compound, and straight-line speed indicators. The visual analysis does not claim to isolate hidden technical parameters of the cars, but it provides a structured and data-driven interpretation of where and under which conditions the main performance differences emerge.
	\end{quote}
	
	% ------------------------------------------------------------
	\section{Final Decision}
	
	The final visualization structure is:
	
	\begin{enumerate}[leftmargin=1.2cm]
		\item \textbf{Competitive Balance Overview}
		
		This dashboard answers what changed between 2021 and 2022.
		
		\item \textbf{Qualifying vs Race Performance}
		
		This dashboard explains whether the change depends on the session context.
		
		\item \textbf{Circuit and Sector DNA}
		
		This dashboard explains where the performance advantage appears.
		
		\item \textbf{Tyre and Straight-Line Performance}
		
		This dashboard adds tyre behaviour and speed indicators, while explicitly declaring the limits of interpretation.
	\end{enumerate}
	
	This structure is the best compromise between analytical depth, Formula 1 storytelling, visual impact, and the five-minute time constraint.
	
\end{document}